\section{Cote 136: the associative localization is terminal}
\label{sec:136associativeterminal}

\paragraph{ARCH-136-ASSOCIATIVE-TERMINAL.}
Let $k$ be any commutative ring. Define $\mathcal B_k$ by
replacing the strict commutative PD algebras in
$\mathcal A_k$ by unital associative differential algebras
with a central coefficient map $S_k=\Gamma_k(T)\to C$,
and retaining both gradings and precisely the same uniform
vertical and horizontal admissibility conditions. The
groupoid obtained from $\mathcal B_k$ by inverting every
arrow is terminal. In particular every component or
localization loop which might remain in the original
integral PD category is invisible after this forgetful
functor. This asserts neither that the forgetful functor
is full nor that it reflects connectedness or equality
of localized arrows.

This concerns exactly the distinction of algebraic
structures in cote 136, physical p.~121, footer 120,
\texttt{ARCHIVE\_136\_scan121.tex}; the retained admissibility
conditions are on physical pp.~117 and 119, footers
116 and 118. It is an extension of the explicit
ARCH-136-ASSOCIATIVE-COMPARISON construction, not a
claim that its intermediate objects belong to
Grothendieck's strict commutative divided-power category.

\begin{proof}
All the constructions above work over $k$, replacing
$\mathbb Z[X]$ by $k[X]$. We verify that their valuewise
zigzag is represented by actual objects of $\mathcal B_k$.
Let
\[
 U:s\operatorname{Mod}k\longrightarrow
                  (\operatorname{Mod}k)^{\mathbb N}
\]
be the degreewise forgetful functor, and let $G$ be its
right adjoint. Explicitly,
\[
 G((A_l)_l)_n=
       \prod_{l\geq0}\prod_{\theta:[l]\to[n]} A_l.
                                                               \tag{136.AT1}
\]
Put $R=GU$. Products in (136.AT1) show that $R$ acts
also on simplicial associative algebras, and that its
unit and multiplication are algebra maps: the unit
evaluates all simplicial restrictions of an element,
and the adjunction counit selects the identity coordinate.
Consequently $C\to RC\rightrightarrows R^2C\cdots$ is
the usual augmented cosimplicial resolution by associative
algebras. Its cochain totalization, with the product in
(136.AC4), is the simplicial associative algebra
$\mathsf K(C)_n=\mathsf K_C(\Delta[n])$ of the preceding
construction. Equivalently its form-degree $p$ row is
\[
          \mathsf K(C)^{w,p}_*
          =\bigoplus_{r+q=p}R^{r+1}(C^{w,q}_*).
                                                               \tag{136.AT2}
\]
To see the agreement with the representable resolution,
write $L=FU$ for the left free resolution functor on
simplicial modules. The adjunctions $F\dashv U\dashv G$
give $L\dashv R$, and linearizing the set resolution
gives $k[Q^{r+1}X]=L^{r+1}k[X]$. Thus
$\operatorname{Hom}(k[Q^{r+1}X],C)=
\operatorname{Hom}(k[X],R^{r+1}C)$, compatibly with
faces, cofaces and the displayed multiplications.

Every $R(M)$ is simplicially contractible, for an
arbitrary simplicial $k$-module $M$. Here are full
module details. Let $G_l(A)$ be the right adjoint
to evaluation in simplicial degree $l$. The natural
Dold--Kan decomposition is
\[
 M_l=\bigoplus_{[l]\twoheadrightarrow[j]}(NM)_j.
\]
In nonnegative homological complexes, the right adjoint
to evaluation at $j$ is the disk $D_{j+1}(A)$, which
has $A\xrightarrow{1}A$ in degrees $j+1,j$:
any map $K_j\to A$ uniquely extends to a chain map
$K\to D_{j+1}(A)$ by the formula in degree $j+1$
given by composition with $\partial:K_{j+1}\to K_j$.
The adjunction and the decomposition therefore give
\[
              NG_l(A)\simeq
                 \prod_{[l]\twoheadrightarrow[j]}
                                      D_{j+1}(A).       \tag{136.AT3}
\]
This is a finite product of contractible disks.
Also $G((A_l))=\prod_l G_l(A_l)$. Normalization
commutes with products; coordinatewise disk contractions
contract the resulting product complex. This proves
the asserted contractibility of $R(M)$ without any
projectivity or injectivity assumption on its entries.
Equation (136.AT2) is a finite sum of such rows, so
every horizontal row of $\mathsf K(C)$ is contractible
in \emph{every} weight and form degree.

For vertical cohomology, the augmented cosimplicial
resolution becomes split after applying $U$, by the
adjunction counit. Its extra codegeneracy gives a
contraction in the resolution direction in each
original simplex degree. This contraction commutes
with the internal differential of $C$. With the total
signs of (136.AC3), the vertical cross terms in the
homotopy equation cancel. Hence
\[
                C_n^{w,*}\longrightarrow
                         \mathsf K(C)_n^{w,*}
\]
is a chain-homotopy equivalence for every $n,w$.
Alternatively, evaluation at $\Delta[n]$ is exact,
so the projective resolution already computes precisely
$C_n^{w,*}$ with no derived correction. It follows
that the uniform vertical bounds of $\mathsf K(C)$
are exactly those of $C$. Thus $\mathsf K$ is an
endofunctor of $\mathcal B_k$.

Apply the same construction to the constant simplicial
algebra $\underline S_k$ in form degree zero, even
though this input is not admissible. The resulting
$\mathsf S_k=\mathsf K(\underline S_k)$ has vertical
cohomology $S_k$ in degree zero in every simplex and
weight, and its every horizontal row is contractible
by (136.AT2). Thus $\mathsf S_k\in\mathcal B_k$.
The augmentation and coefficient maps give natural
transformations of functors on $\mathcal B_k$,
\[
                 \operatorname{Id}\longrightarrow\mathsf K
                          \longleftarrow\operatorname{const}_{\mathsf S_k}.
                                                               \tag{136.AT4}
\]
After all arrows are inverted, both transformations
are natural isomorphisms. For any localized arrow
$f:C\to D$, their resulting natural isomorphism
$e:\operatorname{Id}\simeq\operatorname{const}_{\mathsf S_k}$
satisfies $e_Df=e_C$. It forces every hom-set to
contain at most one arrow; the same isomorphism
supplies an arrow between any two objects. The
localization is nonempty and terminal, as claimed.
\end{proof}
